% !PRINTED 2
\weberpage{2}
%<<CONT
\gesperrt{zahl}. Die Zahl $1$ selbst pflegt man aus
Zweckmässigkeitsgründen nicht als Primzahl zu bezeichnen. Sind zwei
Zahlen durch eine dritte theilbar, so ist auch die Summe und die
Differenz der beiden ersten durch die dritte theilbar; und ist eine
Zahl durch eine zweite, diese durch eine dritte theilbar, so ist auch
die erste durch die dritte theilbar.

Zwei Zahlen haben immer den gemeinsamen Theiler $1$. Wenn sie keinen
anderen gemeinsamen Theiler haben, wie z.~B. die Zahlenpaare $5$ und
$7$ oder $21$ und $38$, so heissen die beiden Zahlen \gesperrt{relative
Primzahlen} oder auch \gesperrt{theilerfremde Zahlen}.

Unter den gemeinsamen Theilern von irgend zwei gegebenen Zahlen wird
einer der grösste sein und es ist eine sehr wichtige Aufgabe, diesen
grössten gemeinschaftlichen Theiler zu finden. Dazu führt ein Verfahren,
das unter dem Namen \gesperrt{Algorithmus des grössten gemeinschaftlichen
Theilers} bekannt ist und sich schon bei Euklid findet.\footnote{Elemente,
Buch VII, II, Bd. II der Heiberg'schen Ausgabe.}

Sind $a,a_1$ die beiden Zahlen, so nehme man die grössere von ihnen,
die $a$ sei, als Dividenden, die kleinere $a_1$ als Divisor, und
bestimme den Quotienten $q_1$; und wenn die Division nicht aufgeht,
den Rest $a_2$, also
\[
 a=q_1a_1+a_2,
\]
so dass $a_2<a_1$ ist. Jeder gemeinschaftliche Theiler von $a$ und
$a_1$ ist dann auch gemeinschaftlicher Theiler von $a_1$ und $a_2$
und umgekehrt. Verfährt man mit $a_1,a_2$ ebenso wie mit $a$ und $a_1$
und setzt, wenn die Division nicht aufgeht,
\[
 a_1=q_2a_2+a_3,
\]
so ist wieder $a_3<a_2$, und jeder gemeinschaftliche Theiler von $a_1$
und $a_2$ ist auch gemeinschaftlicher Theiler von $a_2$ und $a_3$ und
umgekehrt. Fährt man auf diese Weise fort zu dividiren, so muss, da die
Zahlen $a,a_1,a_2,a_3,\ldots$ immer abnehmen, nothwendiger Weise die
Division nach einer endlichen Anzahl von Schritten aufgehen, und es
muss also zuletzt ein Paar von Gleichungen
\[
 a_{\nu-2}=q_{\nu-1}a_{\nu-1}+a_\nu,
 \qquad
 a_{\nu-1}=q_\nu a_\nu
\]
auftreten. Dann ist $a_\nu$ ein gemeinschaftlicher Theiler aller
vorausgehenden $a$, also auch von $a$ und $a_1$, und jeder
gemeinschaftliche Theiler von $a$ und $a_1$ ist Theiler von $a_\nu$.
Also ist $a_\nu$ der grösste gemeinschaftliche Theiler von $a$ und
$a_1$, und wir sind zugleich zu dem Satze gelangt, dass jeder
gemeinschaftliche Theiler zweier Zahlen in ihrem grössten
gemeinschaftlichen Theiler aufgehen muss.
